\begin{thebibliography}{51}
\bibitem{ref1} Ian Agol. The virtual Haken conjecture. Documenta Mathematica, 18:1045–1087, 2013. With an appendix by Ian Agol, Daniel Groves, and Jason Manning.
\bibitem{ref2} Sanjeev Arora and Boaz Barak. Computational Complexity - A Modern Approach. Cambridge University Press, 2009. doi:10.1017/CBO9780511804090.
\bibitem{ref3} David A. Mix Barrington. Bounded-width polynomial-size branching programs recognize exactly those languages in $\mathrm{NC}^1$ . Journal of Computer and System Sciences, 38(1):150–164, 1989. doi:10.1016/0022-0000(89)90037-8.
\bibitem{ref4} David A. Mix Barrington and Denis Thérien. Finite monoids and the fine structure of $\mathrm{NC}^1$ . Journal of the ACM, 35:941–952, 1988. doi:10.1145/48014.63138.
\bibitem{ref5} Laurent Bartholdi, Michael Figelius, Markus Lohrey, and Armin Weiß. Groups with ALOGTIME-hard word problems and PSPACE-complete compressed word problems. Technical report, arXiv.org, 2020. arXiv:1909.13781.
\bibitem{ref6} Laurent Bartholdi, Rostislav I. Grigorchuk, and Zoran Šuniḱ. Branch groups. In Handbook of algebra, Volume 3, pages 989–1112. Elsevier/North-Holland, Amsterdam, 2003. doi: 10.1016/S1570-7954(03)80078-5.
\bibitem{ref7} Laurent Bartholdi and Volodymyr V. Nekrashevych. Iterated monodromy groups of quadratic polynomials. I. Groups, Geometry, and Dynamics, 2(3):309–336, 2008. doi:10.4171/GGD/42.
\bibitem{ref8} Martin Beaudry, Pierre McKenzie, Pierre Péladeau, and Denis Thérien. Finite monoids: From word to circuit evaluation. SIAM Journal on Computing, 26(1):138–152, 1997. doi: 10.1137/S0097539793249530.
\bibitem{ref9} Richard Beigel and John Gill. Counting classes: Thresholds, parity, mods, and fewness. Theoretical Computer Science, 103(1):3–23, 1992. doi:10.1016/0304-3975(92)90084-S.
\bibitem{ref10} William W. Boone. The Word Problem. Annals of Mathematics, 70(2):207–265, 1959. URL: www.jstor.org/stable/1970103.
\bibitem{ref11} Daniel P. Bovet, Pierluigi Crescenzi, and Riccardo Silvestri. A uniform approach to define complexity classes. Theoretical Computer Science, 104(2):263–283, 1992. doi:10.1016/ 0304-3975(92)90125-Y.
\bibitem{ref12} John W. Cannon, William J. Floyd, and Walter R. Parry. Introductory notes on Richard Thompson’s groups. L’Enseignement Mathématique, 42(3):215–256, 1996.
\bibitem{ref13} Moses Charikar, Eric Lehman, Ding Liu, Rina Panigrahy, Manoj Prabhakaran, Amit Sahai, and Abhi Shelat. The smallest grammar problem. IEEE Transactions on Information Theory, 51(7):2554–2576, 2005. doi:10.1109/TIT.2005.850116.
\bibitem{ref14} Max Dehn. Über unendliche diskontinuierliche Gruppen. Mathematische Annalen, 71(1):116– 144, 1911. doi:10.1007/BF01456932.
\bibitem{ref15} Michael R. Garey and David S. Johnson. Computers and Intractability: A Guide to the Theory of NP–completeness. Freeman, 1979.
\bibitem{ref16} Max Garzon and Yechezkel Zalcstein. The complexity of Grigorchuk groups with application to cryptography. Theoretical Computer Science, 88(1):83–98, 1991. doi:10.1016/0304-3975(91) 90074-C.
\bibitem{ref17} Rostislav I. Grigorchuk. Burnside’s problem on periodic groups. Functional Analysis and Its Applications, 14:41–43, 1980. doi:10.1007/BF01078416.
\bibitem{ref18} Rostislav I. Grigorchuk and Zoran Šuniḱ. Asymptotic aspects of Schreier graphs and Hanoi Towers groups. C. R. Math. Acad. Sci. Paris, 342(8):545–550, 2006. doi:10.1016/j.crma. 2006.02.001.
\bibitem{ref19} Victor S. Guba and Mark V. Sapir. On subgroups of the R. Thompson group F and other diagram groups. Matematicheskii Sbornik, 190(8):3–60, 1999. doi:10.1070/ SM1999v190n08ABEH000419.
\bibitem{ref20} Narain Gupta and Saïd Sidki. On the Burnside problem for periodic groups. Mathematische Zeitschrift, 182(3):385–388, 1983. doi:10.1007/BF01179757.
\bibitem{ref21} Frédéric Haglund and Daniel T. Wise. Coxeter groups are virtually special. Advances in Mathematics, 224(5):1890–1903, 2010. doi:10.1016/j.aim.2010.01.011.
\bibitem{ref22} Ulrich Hertrampf. Über Komplexitätsklassen, die mit Hilfe von k-wertigen Funktionen definiert werden. Habilitationsschrift, Universität Würzburg, 1994.
\bibitem{ref23} Ulrich Hertrampf. The shapes of trees. In Proceedings of COCOON 1997, volume 1276 of Lecture Notes in Computer Science, pages 412–421. Springer, 1997. doi:10.1007/BFb0045108.
\bibitem{ref24} Ulrich Hertrampf. Algebraic acceptance mechanisms for polynomial time machines. SIGACT News, 31(2):22–33, 2000. doi:10.1145/348210.348215.
\bibitem{ref25} Ulrich Hertrampf, Clemens Lautemann, Thomas Schwentick, Heribert Vollmer, and Klaus W. Wagner. On the power of polynomial time bit-reductions. In Proceedings of the Eighth Annual Structure in Complexity Theory Conference, pages 200–207. IEEE Computer Society Press, 1993. doi:10.1109/SCT.1993.336526.
\bibitem{ref26} Ulrich Hertrampf, Heribert Vollmer, and Klaus Wagner. On balanced versus unbalanced computation trees. Mathematical Systems Theory, 29(4):411–421, 1996. doi:10.1007/BF01192696.
\bibitem{ref27} Derek F. Holt, Markus Lohrey, and Saul Schleimer. Compressed decision problems in hyperbolic groups. In Proceedings of STACS 2019, volume 126 of LIPIcs, pages 37:1–37:16. Schloss Dagstuhl - Leibniz-Zentrum für Informatik, 2019. URL: http://www.dagstuhl.de/dagpub/ 978-3-95977-100-9.
\bibitem{ref28} Derek F. Holt, Sarah Rees, and Claas E. Röver. Groups, Languages and Automata, volume 88 of London Mathematical Society Student Texts. Cambridge University Press, 2017. doi: 10.1017/9781316588246.
\bibitem{ref29} Birgit Jenner, Pierre McKenzie, and Denis Thérien. Logspace and logtime leaf languages. Information and Computation, 129(1):21–33, 1996. doi:10.1006/inco.1996.0071.
\bibitem{ref30} Howard J. Karloff and Walter L. Ruzzo. The iterated mod problem. Information and Computation, 80(3):193–204, 1989. doi:10.1016/0890-5401(89)90008-4.
\bibitem{ref31} Daniel König and Markus Lohrey. Evaluation of circuits over nilpotent and polycyclic groups. Algorithmica, 80(5):1459–1492, 2018. doi:10.1007/s00453-017-0343-z.
\bibitem{ref32} Daniel König and Markus Lohrey. Parallel identity testing for skew circuits with big powers and applications. International Journal of Algebra and Computation, 28(6):979–1004, 2018. doi:10.1142/S0218196718500431.
\bibitem{ref33} Jörg Lehnert and Pascal Schweitzer. The co-word problem for the Higman-Thompson group is context-free. Bulletin of the London Mathematical Society, 39(2):235–241, February 2007. doi:10.1112/blms/bdl043.
\bibitem{ref34} Martin W. Liebeck, Eamonn A. O’Brien, Aner Shalev, and Pham Huu Tiep. The Ore conjecture. Journal of the European Mathematical Society, 12(4):939–1008, 2010. doi:10.4171/JEMS/220.
\bibitem{ref35} Yury Lifshits and Markus Lohrey. Querying and embedding compressed texts. In Proceedings of MFCS 2006, volume 4162 of Lecture Notes in Computer Science, pages 681–692. Springer, 2006. doi:10.1007/11821069\_59.
\bibitem{ref36} Richard J. Lipton and Yechezkel Zalcstein. Word problems solvable in logspace. Journal of the Association for Computing Machinery, 24(3):522–526, 1977. doi:10.1145/322017.322031.
\bibitem{ref37} Markus Lohrey. Leaf languages and string compression. Information and Computation, 209(6):951–965, 2011. doi:10.1016/j.ic.2011.01.009.
\bibitem{ref38} Markus Lohrey. The Compressed Word Problem for Groups. Springer Briefs in Mathematics. Springer, 2014. doi:10.1007/978-1-4939-0748-9.
\bibitem{ref39} Volodymyr Nekrashevych. Self-similar groups, volume 117 of Mathematical Surveys and Monographs. American Mathematical Society, Providence, RI, 2005. doi:10.1090/surv/117.
\bibitem{ref40} Piotr S. Novikov. On the algorithmic unsolvability of the word problem in group theory. Trudy Mat. Inst. Steklov, pages 1–143, 1955. In Russian. URL: http://mi.mathnet.ru/eng/tm1180.
\bibitem{ref41} David Robinson. Parallel Algorithms for Group Word Problems. PhD thesis, University of California, San Diego, 1993.
\bibitem{ref42} Joseph J. Rotman. An Introduction to the Theory of Groups (fourth edition). Springer, 1995.
\bibitem{ref43} Amir Shpilka and Amir Yehudayoff. Arithmetic circuits: A survey of recent results and open questions. Foundations and Trends in Theoretical Computer Science, 5(3-4):207–388, 2010. doi:10.1561/0400000039.
\bibitem{ref44} Hans-Ulrich Simon. Word problems for groups and contextfree recognition. In Proceedings of FCT 1979, pages 417–422. Akademie-Verlag, 1979.
\bibitem{ref45} Roman Smolensky. Algebraic methods in the theory of lower bounds for boolean circuit complexity. In Proceedings of STOC 1987, pages 77–82. ACM, 1987. doi:10.1145/28395. 28404.
\bibitem{ref46} Jacques Tits. Free subgroups in linear groups. Journal of Algebra, 20(2):250–270, 1972. doi:10.1016/0021-8693(72)90058-0.
\bibitem{ref47} Heribert Vollmer. Introduction to Circuit Complexity. Springer, Berlin, 1999. doi:10.1007/ 978-3-662-03927-4.
\bibitem{ref48} Stephan Waack. The parallel complexity of some constructions in combinatorial group theory. Journal of Information Processing and Cybernetics EIK, 26:265–281, 1990. doi: 10.1007/BFb0029647.
\bibitem{ref49} Jan Philipp Wächter and Armin Weiß. An automaton group with pspace-complete word problem. In Proceedings of STACS 2020, volume 154 of LIPIcs, pages 6:1–6:17. Schloss Dagstuhl - Leibniz-Zentrum für Informatik, 2020. URL: https://www.dagstuhl.de/dagpub/ 978-3-95977-140-5.
\bibitem{ref50} John S. Wilson. Embedding theorems for residually finite groups. Mathematische Zeitschrift, 174(2):149–157, 1980. doi:10.1007/BF01293535.
\bibitem{ref51} Daniel T. Wise. Research announcement: the structure of groups with a quasiconvex hierarchy. Electronic Research Announcements in Mathematical Sciences, 16:44–55, 2009. URL: http: //aimsciences.org//article/id/8d3fc128-8d02-4fee-af67-38001ca1d0ac.
\end{thebibliography}
